% =============================================================================
% §3 Candidate mechanisms (M2) — three separable hypotheses; why smooth models
% struggle with fast collapse; the ABM as the device that separates predictions.
% =============================================================================
\section{Candidate Mechanisms and Their Signatures}
\label{sec:mechanisms}

Three families of mechanism could generate a marriage-composition collapse of
the observed speed. They are not mutually exclusive, but they carry different
observable signatures, which is what makes discrimination possible.

\paragraph{H$_{\text{cohort}}$: cohort replacement.}
Successive birth cohorts enter adulthood with durably lower marriage propensity
--- formed by education expansion, structural change, and the ideational change
of the Second Demographic Transition \citep{callesvogl2026} --- and the stock
composition shifts as these
cohorts flow through the age distribution. The signature: change localizes on
\emph{cohort} lines in an age--period--cohort decomposition; within a cohort,
behavior is comparatively stable; the aggregate can nonetheless accelerate
purely mechanically, as steep-gradient cohorts come to dominate the observed
age range.

\paragraph{H$_{\text{shock}}$: exogenous period shock.}
A calendar-time force common to all living cohorts --- an economic shock, a
policy change, a pandemic --- shifts marriage entry for everyone at once. The
signature: change localizes on \emph{period} lines, visible as parallel shifts
across cohorts at the same calendar date, with no dependence on the prevailing
composition itself.

\paragraph{H$_{\text{cascade}}$: within-cohort reflexive cascade.}
The mechanism behind most ``fast-dynamics'' narratives: individuals respond to
the behavior of their reference group, so that as non-marriage spreads, the
social cost of not marrying falls, generating self-reinforcing acceleration ---
a norm cascade or tipping dynamic. The signature is \emph{state dependence with
an amplifying sign}: within a cohort, marriage decline must be \emph{larger}
where the reference-group not-married share is already high. This is the only
mechanism of the three that is genuinely reflexive, and identifying it from
observational data confronts the classical reflection problem
\citep{manski1993}; Section~\ref{sec:discrimination} states how our design
addresses the sign-based, within-cohort version of the test.

Why not adjudicate with a standard smooth model? Because the phenomenon's
defining feature --- late acceleration --- is exactly what smooth aggregate
specifications absorb rather than explain. A logistic trend fitted to the
married share will always ``fit'' a collapse after the fact; it cannot say
whether the curvature came from cohort turnover, a dated shock, or feedback.
We therefore proceeded in two deliberately separated steps. First, the
econometric discrimination of Section~\ref{sec:discrimination}, which asks the
data directly which signature is present. Second, a generative test: we built
an agent-based implementation of the \emph{cascade} hypothesis --- a
coordination threshold in partnership formation, the formalization of the
tipping narrative --- and calibrated it to the composition data. That model was
a genuine, live candidate, not a straw man: with its threshold active it
generated a real, endogenous marriage decline, and removing the threshold
flipped the generated fertility path from collapse to increase (the mechanism
was doing the work). It nonetheless failed in a specific, diagnostic way: the
calibrated threshold produced a smooth glide --- roughly $-0.02$ per year in
generated TFR, about 55 percent of the observed magnitude --- with none of the
observed late acceleration, and its failure shaped the discrimination design
that follows. Cascade dynamics, when actually calibrated to composition data,
do not produce the observed shape; the question became which non-reflexive
structure does.
